There's a zoo of NP-complete problems. For any pair $(L,H)$ of \textbf{NP}-complete problems, we have that:
$$L \leq_{p} H \ H \leq_{p} L$$
meaning that $L$ and $H$ are equivalent. By the mechanism of reductions, we have an easy way to show the equivalances between decision problems, even though
for one of them does not exist a polytime algorithm solving it. Being said this, defining a reduction from $L$ to $H$ sometimes is easy, but sometimes is very 
difficult, and it is useful to think at \textbf{NP}-complete problems as a graph, where edges are natural polytime reductions. \vspace{7pt}

As we already guessed, between \textbf{P} and \textbf{EXP} there is \textbf{NP}, a complexity class from which we can distinguish hard problems from complete 
problems.